\documentclass[hardcover,french,color]{scriptorium}

\title{Le Grand Œuvre.exe}
\author{Calion}

\publicationyear{2026}

\website{libricalionis.com}
\contactemail{libri.calionis@gmail.com}

\titleimagecolor{assets/alchemist_color.png}
\titleimagegrayscale{assets/alchemist_grayscale.png}

\begin{document}

  \maketitle

  \mainmatter

  \chapter{La pierre philosophale nécessite Java 17}

  \section{Le laboratoire}

  \lettrine{B}{althazar}, quarante-deux ans, avait enfin réussi à transformer le plomb en or.

  Le problème était que son script Python refusait de l'admettre.

  \scenebreak

  \scriptoriumnote
  {
    Toute transmutation doit être reproductible avant de pouvoir être qualifiée de Grand Œuvre.
  }

  Balthazar relança le programme.

  \scriptoriumquote
  {
    ERROR 418

    Substance obtenue : Au

    Substance attendue : Pb
  }

  Son téléphone vibra.

  \smsa{Tu viens manger ?}

  \smsb{Dans cinq minutes. J'ai un problème avec l'or.}

  \smsa{Encore ?}

  \calionnote
  {
    Les anciens cherchaient la pierre philosophale.

    Les modernes cherchent pourquoi elle fonctionne sur leur machine mais pas en production.
  }

  Balthazar éteignit l'ordinateur.

  Le lingot redevint du plomb.

  Il le ralluma.

  De l'or.

  Balthazar ouvrit un ticket.

  \theend

\end{document}
